% META: {"id": "TEXP-CTP-CO-040", "chapitre": "TEXP-COMPLEXES-TRIGO-POLYNOMES", "type_objet": "corrige", "exercice_ref": "TEXP-CTP-EX-040", "capacites_codes": ["C5"], "sources_inspiration": [], "status": "approved"}
% BEGIN-VERIFY
% from sympy import *
% x = symbols('x')
% # f(x) = sqrt(1+x), tangente en 0: f(0)=1, f'(0)=1/2, y=1+x/2
% f = sqrt(1 + x)
% fp = diff(f, x)
% assert fp.subs(x, 0) == Rational(1, 2)
% # f concave: f'' < 0
% fpp = diff(f, x, 2)
% assert fpp == -1/(4*(1+x)**Rational(3,2))
% assert fpp.subs(x, 0) == -Rational(1,4)  # concave
% # approximation: sqrt(1.1) ~ 1.05
% approx = 1 + Rational(1,20)  # 1 + 0.1/2
% assert approx == Rational(21, 20)
% exact = f.subs(x, Rational(1,10))
% assert exact == sqrt(Rational(11,10))
% # concave => courbe sous la tangente => approx >= exact
% assert approx > float(exact)
% END-VERIFY
\begin{corrige}{TEXP-CTP-EX-040}

\textbf{Question 1.} $f(0)=1$, $f'(x) = \dfrac{1}{2\sqrt{1+x}}$, $f'(0) = \dfrac{1}{2}$.

\commentaireMarge{$\sqrt{1+x}$ concave ($f''<0$) : la tangente majore la courbe.}

Tangente : $y = 1 + \dfrac{x}{2}$. Pour $x = 0{,}1$ : $\sqrt{1{,}1} \approx 1 + \dfrac{0{,}1}{2} = 1{,}05$.

\textbf{Question 2.} $f''(x) = -\dfrac{1}{4(1+x)^{3/2}} < 0$ : $f$ est concave, donc la courbe est \textbf{au-dessous} de la tangente.

L'approximation $1{,}05$ est donc \textbf{par exces}.

\end{corrige}
